\chapter{Software Requirements Specification}
\label{app:srs}

This appendix contains the complete Software Requirements Specification (SRS) for the Dynamic \qos Scheduler system.

\section{Functional Requirements}

\subsection{VPN Tunnel Management}

\textbf{FR-01} The system shall establish a local VpnService tunnel on the host device when the user activates QoS scheduling.

\textbf{FR-02} The system shall route all hotspot traffic through the TUN interface created by the VpnService.

\textbf{FR-03} The system shall forward processed packets to the internet uplink after classification and scheduling.

\textbf{FR-04} The system shall gracefully tear down the VPN tunnel when the user deactivates QoS scheduling or the hotspot is disabled.

\textbf{FR-05} The system shall automatically restart the tunnel if it drops unexpectedly while the hotspot remains active.

\subsection{Device Discovery and Registry}

\textbf{FR-06} The system shall detect all devices connected to the hotspot by inspecting source IP addresses from intercepted packets.

\textbf{FR-07} The system shall maintain a registry of connected devices including: IP address, MAC address (where resolvable), hostname (where resolvable), assigned priority class, and per-session traffic statistics.

\textbf{FR-08} The system shall remove a device from the active registry after a configurable inactivity timeout (default: 60 seconds).

\textbf{FR-09} The system shall persist user-assigned priority overrides across sessions using device MAC address as the key.

\subsection{Traffic Classification (DPI-Lite)}

\textbf{FR-10} The system shall parse raw IPv4 and IPv6 packet headers from the TUN interface using Java NIO ByteBuffer.

\textbf{FR-11} The system shall extract the 5-tuple from each packet: source IP, destination IP, source port, destination port, protocol (TCP/UDP).

\textbf{FR-12} The system shall classify each flow into application categories based on destination port and protocol.

\textbf{FR-13} The system shall allow the user to manually override the traffic class of any active flow or device.

\textbf{FR-14} The system shall re-classify flows periodically (every 5 seconds) to adapt to changing traffic patterns.

\subsection{Bandwidth Scheduling (Token Bucket)}

\textbf{FR-15} The system shall instantiate one token bucket per connected device.

\textbf{FR-16} Each token bucket shall be configurable with rate (bytes/second) and burst (bytes) parameters.

\textbf{FR-17} The system shall map priority classes to default token bucket parameters (HIGH=80\%, MEDIUM=50\%, LOW=20\% of uplink).

\textbf{FR-18} The system shall drop or delay packets that exceed the token bucket allowance for their device.

\textbf{FR-19} The system shall dynamically adjust token bucket rates when devices join or leave the hotspot.

\textbf{FR-20} The system shall allow the user to set a manual bandwidth cap per device.

\section{Non-Functional Requirements}

\subsection{Performance}

\textbf{NFR-01} Packet processing latency shall not exceed 5 ms per packet under normal load.

\textbf{NFR-02} The scheduler shall sustain processing of at least 10,000 packets/second.

\textbf{NFR-03} CPU utilization shall remain below 15\% on mid-range Android devices.

\textbf{NFR-04} Memory footprint shall remain below 80 MB RAM.

\subsection{Reliability}

\textbf{NFR-05} The VPN tunnel shall recover automatically within 3 seconds of disconnection.

\textbf{NFR-06} The system shall not cause packet loss beyond the configured token bucket drop policy.

\textbf{NFR-07} Priority assignments shall survive application restarts and device reboots.

\subsection{Security}

\textbf{NFR-10} The system shall not log or persist packet payload content.

\textbf{NFR-11} The VPN tunnel shall be local-only (no external routing).

\textbf{NFR-12} Persisted device priority data shall be stored in Android internal storage.

\subsection{Compatibility}

\textbf{NFR-13} The application shall support Android API level 29 (Android 10) and above.

\textbf{NFR-14} The application shall function correctly on both IPv4 and IPv6 hotspot configurations.

\textbf{NFR-15} The application shall not require root access, ADB access, or kernel modifications.

\section{Acceptance Criteria}

\subsection{AC-01: Tunnel Establishment}
\textbf{Given} the Android hotspot is active and the user has granted VPN permission\\
\textbf{When} the user taps "Start QoS" in the app\\
\textbf{Then} a VpnService tunnel is established within 2 seconds

\subsection{AC-15: HIGH Priority Throughput Advantage}
\textbf{Given} two devices are connected — Device A assigned HIGH, Device B assigned LOW\\
\textbf{When} both devices simultaneously run iPerf3 to saturate the uplink\\
\textbf{Then} Device A achieves at least 3× the throughput of Device B

\subsection{AC-26: Packet Processing Latency}
\textbf{Given} 5 devices are connected with 100 Mbps aggregate throughput\\
\textbf{When} per-packet processing time is measured\\
\textbf{Then} the 99th-percentile processing latency is at or below 5 ms

\textit{(Complete SRS with all 32 acceptance criteria available in project documentation)}
